\documentclass[a4paper,10.5pt]{article}
\usepackage[T1]{fontenc}\usepackage[utf8]{inputenc}
\usepackage[italian]{babel}
\usepackage[top=2.2cm,bottom=2.2cm,left=2.35cm,right=2.35cm]{geometry}
\usepackage{hyperref}\usepackage{xcolor}\definecolor{accent}{HTML}{1F4E78}\hypersetup{colorlinks=true,urlcolor=accent}
\pagestyle{empty}\setlength{\parindent}{0pt}\setlength{\parskip}{7pt}
\begin{document}
\textbf{Prabhavit Gupta}\\
Venezia, Italia \textbar{} +39 351 610 7807 \textbar{} prabhavitg@gmail.com\\[10pt]
Gentile Team di Lutech,

vi scrivo per candidarmi alla posizione di ML Engineer per la Practice Data Analytics. Negli ultimi tre anni ho lavorato a livello di conoscenza pratica con pipeline di Machine Learning e AI, con particolare attenzione a modelli generativi (LLM) e sistemi agentic in produzione.

In FlowSpace ho costruito e mantengo un bot AI che genera annunci di spazi di lavoro a partire da brochure in qualsiasi formato (PDF, sito web, foto o foglio di calcolo), con un umano nel ciclo di revisione, oltre 900 bozze prodotte finora. Ho anche costruito agenti di supporto AI che rispondono da dati live della piattaforma e passano il caso a una persona quando la confidenza della risposta è bassa, un meccanismo di validazione simile a quello che chiedete per il monitoraggio dei modelli in produzione.

In UltraStrategy ho costruito la pipeline di parsing dei CV integrata nella piattaforma con n8n ed Eden AI: l'utente carica un CV, la pipeline estrae i dati e valuta la rilevanza del candidato per specifici tipi di ruolo, e infine viene generato un CV standardizzato tramite le API di Claude. Ho lavorato quotidianamente con Python per script di automazione e integrazione, e ho seguito da vicino la coerenza dei risultati del modello nel tempo, un'attenzione diretta a KPI tecnici e di business.

Sono onesto sul livello della mia esperienza: il mio lavoro con AI/ML è a livello di integrazione, automazione e validazione pratica di pipeline LLM e RAG, non di ricerca o fine-tuning avanzato di modelli, e non ho ancora esperienza diretta con Docker/Kubernetes in produzione, anche se imparo rapidamente nuovi strumenti.

Vivo a Venezia, ho un permesso di soggiorno UE di lungo periodo e non necessito di sponsorizzazione per lavorare in Italia. Sono disponibile per la modalità ibrida a Bari o Roma e per un colloquio quando preferite. Il mio italiano è conversazionale (B1-B2) e il mio inglese è C2, superiore al livello B2 richiesto.

Grazie per l'attenzione.

Cordiali saluti,\\
Prabhavit Gupta
\end{document}
